\documentclass[12pt,a4paper]{article}

% === Packages ===
\usepackage[utf8]{inputenc}
\usepackage[T1]{fontenc}
\usepackage{amsmath,amssymb,amsthm}
\usepackage{graphicx}
\usepackage{booktabs}
\usepackage{multirow}
\usepackage{natbib}
\usepackage{hyperref}
\usepackage{ragged2e}
\Urlmuskip=0mu plus 1mu\relax
\usepackage[margin=2.5cm]{geometry}
\usepackage{setspace}
\onehalfspacing
\graphicspath{{figures/}}

% Resolve cross-references to the online appendix
\usepackage{xr-hyper}
\externaldocument{online_appendix}

% Public replication package and online appendix
\newcommand{\onlineappendixurl}{https://github.com/vr-rodrigues/jornada-6x1-replication}

% === Metadata ===
\title{How much productivity is needed to shorten the workweek?\\
Brazil's 36-hour proposal and 40-hour alternatives}
\author{Victor Rangel\thanks{Insper. E-mail: \texttt{victorr@insper.edu.br}.}}
\date{April 2026}

\begin{document}

\maketitle

% === Abstract ===
\begin{abstract}
Brazil's workweek debate began with a 36-hour proposal, but it now includes 40-hour and phased alternatives. This paper computes $A_{\text{req}}$, the one-off TFP gain that would make shorter formal hours output-neutral in a short-run accounting model with informality, firm-size heterogeneity, and an hours-productivity schedule. A 40-hour cap requires a modest productivity gain and leaves the representative agent GHH diagnostic positive. The 36-hour endpoint is more demanding, with an output-neutrality benchmark around $6$--$8\%$ depending on the hours-efficiency specification and disciplined calibration envelopes of roughly $5.6$--$7.5\%$ and $6.9$--$9.2\%$. The benchmark is large relative to Brazil's recent TFP record, while most of the output cost comes from fewer formal hours rather than from reallocation into informality. The results support evaluating intermediate and phased paths separately from a direct move to 36 hours; they do not rank reforms by social welfare.
\end{abstract}

\textbf{Keywords:} working hours, informality, total factor productivity, CES aggregator, Brazil, welfare

\textbf{JEL:} E24, J22, J46, O17, O47

\clearpage

% ==================================================================
% 1. Introduction
% ==================================================================
% =============================================================
% Section: Introduction
% =============================================================

\section{Introduction}\label{sec:intro}

Brazil is debating an eight-hour cut to its standard workweek. PEC 8/2025 would lower the cap from 44 to 36 hours and bind roughly two thirds of formal employees at once. Two numbers anchor the policy discussion. The first is the output removed on impact. The second is the one-off productivity gain required to restore it.

The Brazilian setting complicates both answers. Informality is $37.8\%$ under the narrow PNAD definition and $44.2\%$ under the broad IBGE one, ranging from $41\%$ in services to $72\%$ in agriculture. Standard treatments of hours reform model a fully formal economy \citep{Cacciatore2016} or treat informality as a static size distortion \citep{Ulyssea2018}. Neither has been applied to a large cap cut.

This paper builds a static partial-equilibrium model disciplined by Brazilian microdata. A CES aggregator over formal and informal effective labor governs reallocation. A fatigue function $e(h)$ with peak at $40$ hours maps the cap into per-hour productivity. Two firm-size groups generate differential incidence and GHH preferences support welfare. The substitution elasticity is pinned to the Brazilian formal--informal wage premium of \citet{MeghirNaritaRobin2015}, giving $\sigma_{\text{sub}}=1.326$ with interval $[1.116,\,1.469]$ at $(\omega,\eta_I)=(0.622,\,0.40)$.

The central object is $A_{\text{req}}$, the one-off TFP gain that restores pre-reform output. Under the preferred flat-below specification $A_{\text{req}}=6.63\%$. Under the conservative symmetric specification $A_{\text{req}}=8.18\%$. The joint $(\sigma,\omega,\eta_I)$ envelope is $[5.62\%,\,7.48\%]$ under the preferred specification and $[6.93\%,\,9.23\%]$ under the conservative one. At Brazil's best post-1990 decade (PWT \texttt{rtfpna}, $+0.26\%$ per year, 2003--2013) accumulating $A_{\text{req}}$ takes $25$--$31$ years. A stress test that layers a $1$--$2\%$ post-reform TFP contraction raises the conservative $A_{\text{req}}$ to $9.28$--$10.39\%$.

Aggregate informality rises by $+1.6$--$1.9$ pp. With $\sigma>1$ and $\eta_I=0.40$, formal and informal labor are imperfect substitutes, so the cap pushes marginal formal workers toward the informal sector instead of being absorbed by intensive adjustment alone. Welfare peaks near a $41$-hour cap and turns negative below $38$.

The paper makes three contributions. First, it embeds a formal--informal reallocation margin into a structural analysis of hours reform and shows the margin is modest under wage-premium-disciplined parameters. Second, it proposes $A_{\text{req}}$ as a benchmark commensurable with the historical TFP record. Third, it provides a wage-premium-conditional interval for $\sigma_{\text{sub}}$ that the reader can reassess.

Scope is limited. The framework is static and partial equilibrium. It omits capital deepening, general-equilibrium price adjustment, worker selection, and work intensification. Evidence from Portugal 1996 \citep{AsaiLopesTondini2024} and Brazil 1988 \citep{GonzagaMenezesFilhoCamargo2003} enters as a directional plausibility check rather than an identified lower bound for Brazil's effective $A_{\text{req}}$.

Section~\ref{sec:framework} presents facts, model, and calibration. Section~\ref{sec:validation} validates the fit. Section~\ref{sec:results} reports the counterfactual, the TFP benchmark, and the stress test. Section~\ref{sec:conclusion} draws policy implications.


% ==================================================================
% 2. Brazilian facts and framework
% ==================================================================
\section{Brazilian facts and model}\label{sec:framework}

% =============================================================
% Subsection: Brazilian facts (inside Section 2)
% =============================================================

Formal hours are concentrated near the statutory maximum. The baseline maps the DIEESE contracted-hours distribution into three model bins: $8.5\%$ of formal workers in the $\leq 39$h category are assigned to the 36h bin, $26.9\%$ at 40h remain in the 40h bin, and $64.6\%$ in the $\geq 41$h category are assigned to the 44h bin \citep{DIEESE2024}. PNAD Cont\'{i}nua Q4 2024 habitual hours give a flatter distribution, with an average of $41.2$ hours \citep{PNAD2025}. Most formal workers operate at or above 40 hours, the margin a 36-hour cap would bind.

Informality is large and dispersed. The national rate is $37.8\%$ under the narrow PNAD definition and $44.2\%$ under the broad IBGE one \citep{PNAD2025}. The sectoral range runs from $41.2\%$ in services to $71.8\%$ in agriculture (Table~\ref{tab:facts_informality}).

The formal-informal wage premium is positive and moderate. \citet{MeghirNaritaRobin2015} document a conditional wage premium $R\in[1.15, 1.55]$, the ratio of mean formal hourly wages to mean informal hourly wages among otherwise comparable workers. The calibration below uses this premium to pin down the CES substitution elasticity $\sigma_{\text{sub}}$.

Brazilian TFP has been stagnant or declining throughout the post-1990 period. The Penn World Table 11.0 \texttt{rtfpna} series shows an annualized decline of $-0.83\%$ per year from 1990 to 2023 and $-0.82\%$ from 1990 to 2019, before COVID \citep{FeenstraInklaarTimmer2015,PWT2025,FREDPWTBrazil2026}. The best rolling post-1990 decade in the updated series is 2000 to 2010, with only $+0.05\%$ per year. The TFP benchmark in Section~\ref{sec:results} uses this updated series as the anchor for $A_{\text{req}}$.

These facts call for a model with formal workers concentrated near the cap, substitution between formal and informal effective labor, and per-hour productivity that varies with the workweek. The model adds two firm-size groups for incidence and GHH preferences for welfare.

\begin{table}[htbp]
\centering
\caption{Sectoral and national labor-market moments, PNAD Cont\'{i}nua Q4 2024.}\label{tab:facts_informality}
\small
\begin{tabular}{lcccccc}
\hline\hline
Sector & $\lambda_s$ & Informality & Avg.\ hours & $\theta_{36}$ & $\theta_{40}$ & $\theta_{44}$ \\
 & & (broad) & (formal) & & & \\
\hline
Agriculture & 0.076 & 71.8\% & 44.1 & 0.101 & 0.269 & 0.630 \\
Industry    & 0.205 & 44.4\% & 42.5 & 0.048 & 0.417 & 0.535 \\
Services    & 0.720 & 41.2\% & 40.8 & 0.171 & 0.410 & 0.419 \\
\hline
National    & 1.000 & 44.2\% & 41.2 & 0.143 & 0.406 & 0.451 \\
\hline\hline
\end{tabular}
\\[2pt]
\begin{minipage}{0.95\textwidth}
\justifying
{\footnotesize \textit{Notes:} $\lambda_s$ is the sector's share of national employment. Informality uses the broad IBGE definition (VD4009 combined with V4017 CNPJ indicator). ``Avg.\ hours'' is the mean weekly hours of \emph{formal} workers (V4039, habitual). $\theta_{36}$, $\theta_{40}$, $\theta_{44}$ are the shares of formal workers in each contracted-hours bin. \textit{Source:} PNAD Cont\'{i}nua Q4 2024, IBGE microdata, survey-weighted, $N=209{,}311$ \citep{PNAD2025}.}
\end{minipage}
\end{table}

% =============================================================
% Subsection: Model (inside Section 2)
% =============================================================

A representative firm in group $g\in\{\text{Small},\text{Large}\}$ produces output $Y_g$ from capital $K_g$ and effective labor $L_g$ using a Cobb-Douglas technology. Small firms have at most 49 employees, and large firms have 50 or more.
\begin{equation}\label{eq:production}
    Y_g = A\cdot K_g^{\alpha}\cdot L_g^{1-\alpha},
\end{equation}
where $A$ is aggregate productivity and $\alpha$ is the capital share of income. Effective labor $L_g$ combines formal ($F$) and informal ($I$) workers through a CES aggregator,
\begin{equation}\label{eq:ces}
    L_g = \left[\omega L_{F,g}^{\rho} + (1-\omega)L_{I,g}^{\rho}\right]^{1/\rho},\quad \rho=(\sigma_{\text{sub}}-1)/\sigma_{\text{sub}},
\end{equation}
where $\omega$ is the CES weight on formal labor and $\sigma_{\text{sub}}$ is the elasticity of substitution between formal and informal labor. When $\sigma_{\text{sub}}<1$ the two inputs are complements, and raising the cost of formal labor does not trigger a large reallocation toward the informal margin. The two labor inputs are
\[
L_{F,g}=N_{F,g}\sum_b\theta_b h_b e(h_b),\qquad L_{I,g}=\eta_I N_{I,g} h_I e(h_I),
\]
where $N_{F,g}$ and $N_{I,g}$ are the headcounts of formal and informal workers in group $g$ (with $N_{I,g}\equiv N_g-N_{F,g}$ and total group employment $N_g$ fixed), $h_b$ indexes the three contracted-hours bins $b\in\{36,40,44\}$ with shares $\theta_b$ among formal workers, $h_I=44$ is weekly hours for informal workers, and $\eta_I$ is the per-head productivity of an informal worker relative to a formal worker. Per-hour efficiency follows $e(h)=\exp\{-\kappa(h-h^*)^2\}$ above the peak $h^*=40$, where $\kappa$ is the curvature of the long-hours fatigue penalty. Below the peak the paper considers two transparent regimes. The preferred \emph{one-sided fatigue} specification sets $e(h)=1$ for $h\leq h^*$: long workweeks are tiring, but shorter workweeks are not mechanically punished. The conservative \emph{two-sided penalty} specification keeps the same quadratic form below $h^*$: deviations on either side of 40 hours reduce per-hour efficiency. The one-sided specification is informed by \citet{Pencavel2015}, who documents the fatigue leg above the peak but does not characterize efficiency below it. The two-sided specification serves as a conservative upper bound on the required TFP gain.

Each firm chooses the number of formal workers $N_{F,g}$ to maximize profit net of three cost wedges,
\begin{equation}\label{eq:firm}
    Y_g-\tau_g N_{F,g}-\tfrac{1}{2}\pi_m N_{I,g}^2-\tfrac{1}{2}\gamma_{F,g}(N_{F,g}-\bar{N}_{F,g})^2,
\end{equation}
where $\tau_g$ is a linear per-worker cost of formality (CLT/FGTS obligations and size-dependent enforcement), $\pi_m$ is a convex cost of informality (fines, reputational risk, and limited access to formal inputs), and $\gamma_{F,g}$ is a quadratic adjustment cost penalizing deviations from the pre-reform formal headcount $\bar{N}_{F,g}$. The wedge $\tau_g$ is calibrated to match group-specific informality and plays the role of a size-dependent distortion in the spirit of \citet{HsiehKlenow2009} and \citet{RestucciaRogerson2008}, with endogenous-formality microfoundations in \citet{DixCarneiroGoldbergMeghirUlyssea2026}. Cutting the statutory cap from $h_0=44$ hours to $h_1$ reduces effective labor through both the hours bins $h_b$ and the efficiency function $e(h_b)$, and widens the wedge between the marginal product of formal labor and its private cost.

The object of interest is the TFP multiplier that restores aggregate output,
\begin{equation}\label{eq:areq}
    A_{\text{req}}:\quad \sum_g Y_g\!\left(A\,(1+A_{\text{req}}),K_g,h_1\right)=\sum_g Y_g(A,K_g,h_0).
\end{equation}
$A_{\text{req}}$ is the one-off percentage gain in aggregate productivity $A$ that exactly offsets the short-run output loss from moving the cap from $h_0$ to $h_1$. It is an accounting benchmark to compare with Brazil's historical TFP record, not a predicted growth rate.

For welfare, preferences are GHH \citep{GreenwoodHercowitzHuffman1988} with $\sigma_{\text{ghh}}=1$ and Frisch elasticity $1/\nu=0.5$. The compensating variation is
\begin{equation}\label{eq:cv}
    \Delta\text{CV}=\frac{C_1-\psi h_1^{1+\nu}/(1+\nu)}{C_0-\psi h_0^{1+\nu}/(1+\nu)}-1.
\end{equation}
The numerator and denominator are the post- and pre-reform consumption-equivalents of welfare under GHH preferences, computed as consumption net of the monetized labor disutility $\psi h^{1+\nu}/(1+\nu)$. $\Delta$CV is the percentage change in welfare.

The model is static and partial-equilibrium. Capital is predetermined, total employment $N$ is fixed, and general-equilibrium feedbacks (wage-price, monetary, fiscal, cross-sector) are absent. Appendix Table~\ref{tab:omitted_channels} lists the modeled, bracketed, and omitted channels, including compliance, overtime, wage bargaining, labor supply, and sectoral linkages. A textbook user-cost elasticity $\varepsilon_K\in[0.2,0.4]$ \citep{Caballero1999} with capital share $\alpha=0.35$ implies
\begin{equation}\label{eq:pe_bound}
    \Delta A_{\text{req}}^{K}\approx -\alpha\cdot\varepsilon_K\cdot|\Delta L_F/L|\in[-1.0,-1.8]\text{ pp},
\end{equation}
where $\Delta A_{\text{req}}^{K}$ is the downward correction to the short-run $A_{\text{req}}$ once the endogenous response of the capital stock is factored in. The medium-run analog of the headline falls by $1$ to $2$ pp over a $3$ to $5$ year horizon.

% =============================================================
% Subsection: Calibration (inside Section 2)
% =============================================================

Table~\ref{tab:params} lists the parameters under three headings. External parameters come from PWT, DIEESE, and PNAD Cont\'{i}nua. Moment-matched parameters ($\tau_g$, $\pi_m$, $\kappa$, $\psi$) are calibrated internally. Author-chosen parameters ($\eta_I$, $e_q$, $h^*$, $h_I$, $\gamma_F^g$) are set inside ranges consistent with the work-hours evidence of \citet{Pencavel2015} and \citet{CollewetSauermann2017}, the informal-productivity literature \citep{LaPortaShleifer2014}, and the CLT/FGTS compliance-cost literature. Cross-country hours evidence in \citet{BickFuchsSchundelnLagakos2018} provides directional support for declining marginal product of hours but does not identify a peak. The online appendix reports the source and sensitivity for each parameter.

\begin{table}[htbp]
\centering
\caption{Parameter values by evidential basis.}\label{tab:params}
\small
\begin{tabular}{@{}l p{4.2cm} c p{4.8cm}@{}}
\hline\hline
Symbol & Description & Value & Source \\
\hline
\multicolumn{4}{l}{\textit{External}} \\
\hline
$\alpha$     & Capital share                  & 0.35   & PWT 10.01 + \citet{Gollin2002} \\
$\omega$     & CES formal weight              & 0.622  & $1-$ narrow informality (PNAD Cont\'{i}nua 2025) \\
$\theta_b$   & Hours-bin weights (formal)     & $(0.085, 0.269, 0.646)$ & DIEESE 2024 \\
$\nu$        & Inverse Frisch elasticity       & 2.0    & Standard (Frisch $=0.5$) \\
\hline
\multicolumn{4}{l}{\textit{Moment-matched}} \\
\hline
$\sigma_{\text{sub}}$ & CES substitution elasticity & \textbf{1.326} $[1.116,1.469]$ & Wage premium \citep{MeghirNaritaRobin2015} \\
$\tau_S,\tau_L$       & Formal wedges (small, large) & 2.609,\,0.004 & Bisection to group informality \\
$\pi_m^S,\pi_m^L$     & Informality convex cost      & 0.000,\,42.11 & Bisection \\
$\kappa$               & Efficiency curvature         & $2.11\times 10^{-3}$ & Derived from $e_q$ at $h_{\text{ref}}$ \\
$\psi$                 & GHH disutility weight        & $5.74\times 10^{-5}$ & Baseline FOC \\
\hline
\multicolumn{4}{l}{\textit{Author-chosen (within ranges consistent with evidence)}} \\
\hline
$\eta_I$      & Informal relative efficiency   & 0.40   & \citet{LaPortaShleifer2014} informal/formal wage ratio midpoint ($\approx 1/3$ to $1/2$) \\
$e_q$          & Hours--output elasticity       & 0.60   & \citet{Pencavel2015} range \\
$h^*$          & Peak-efficiency hours          & 40     & \citet{Pencavel2015} \\
$h_I$          & Informal hours                 & 44     & No statutory cap in informal sector \\
$\gamma_F^S, \gamma_F^L$ & Adjustment costs        & 0.12,\,0.03 & RAIS + CLT/FGTS \\
\hline\hline
\end{tabular}
\end{table}

The CES substitution elasticity $\sigma_{\text{sub}}$ is disciplined by the wage premium. From the CES first-order condition, the formal--informal wage premium $R$ is a monotone function of $\sigma_{\text{sub}}$. Plugging a widened window around \citeauthor{MeghirNaritaRobin2015}'s estimates, $R\in[1.15,1.55]$, into the CES FOC at the recalibrated $(\omega,\eta_I)=(0.622,0.40)$ yields
\begin{equation}\label{eq:sigma_interval}
R\in[1.15,1.55]\implies \sigma_{\text{sub}}\in[1.116,1.469],\quad A_{\text{req}}\in[7.50\%,8.58\%]\text{ (conservative, midbox).}
\end{equation}
The central value $\sigma=1.326$ corresponds to $R\approx 1.40$, the midpoint of the MNR window.\footnote{The window $R\in[1.15,1.55]$ spans the full range of MNR's reported estimates across sub-populations and alternative specifications (Tables~3 and~10). The mapping from $R$ to $\sigma_{\text{sub}}$ depends on $(\omega,\eta_I)$. The anchors $\omega=0.622$ ($1$ minus the narrow PNAD 2025 informality rate) and $\eta_I=0.40$ (midpoint of the La Porta--Shleifer informal/formal wage ratio range $1/3$ to $1/2$) are theoretically consistent: under competitive wage-taking the CES aggregator identifies $\eta_I$ through the marginal-product (wage) ratio rather than through value-added per employee, which is contaminated by capital-per-worker differences. The online appendix documents the derivation.} Letting $\omega$ and $\eta_I$ also vary over their data-consistent ranges ($\omega\in[0.58,0.66]$, $\eta_I\in[0.33,0.50]$), the joint $(\sigma,\omega,\eta_I)$-box envelope under the conservative specification is $[6.93\%, 9.23\%]$ at the eight corners. Under the preferred flat-below specification the central $A_{\text{req}}$ is $6.63\%$ and the envelope is $[5.62\%,7.48\%]$. The online appendix reports the full grids. The small reallocation elasticity documented by \citet{DerenoncourtGerardLagosMontialoux2025} supports a modest intra-firm $\sigma$ and is distinct from the high substitution elasticities between worker types reported in refugee studies. The online appendix discusses this distinction.

Rebuilding the calibration on PNAD habitual hours ($\theta=(0.143,0.406,0.451)$, $44.2\%$ informality) yields $A_{\text{req}}=7.80\%$ under the conservative specification. The baseline uses DIEESE contracted hours because the statutory cap is a contracted-hours concept.


% ==================================================================
% 3. Calibration checks
% ==================================================================
% =============================================================
% Section: Validation
% =============================================================

\section{Calibration checks and directional benchmarks}\label{sec:validation}

The checks in this section have a narrow role. The model matches the calibration targets by construction. The external benchmarks are directional plausibility checks, not causal validation of the counterfactual. Magnitudes from the Portuguese 1996 reform and the Brazilian 1988 reform should be compared in sign and order of magnitude, not in level.

Panel B1 evaluates the model against domestic moments that are not used to fit the main parameters. The model is close on mean formal hours and on the share of formal workers above $36$h. It understates the cross-sector range in average formal hours, which is a limitation of the parsimonious national hours distribution. As a check, the replication package perturbs the formal hours distribution $\theta$ within a $\pm 5$ pp envelope around the DIEESE 2024 weights and re-runs the calibration. The fitted Panel B1 moments shift by $0.4$ to $0.6$ h/week and $1$ to $2$ pp, with the sign of the gap preserved. These non-targeted moments are informative, but they are not a separate identification exercise.

Panel B2 compares the model with two hours reforms. The Portuguese 1996 reform has a common $44\to 40$h cap, and the Brazilian 1988 reform has a $48\to 44$h cap. Signs match. Aggregate output falls and hourly productivity rises in both data and model. Magnitudes do not match. The firm-level Portuguese estimate of $+4.4\%$ in hourly productivity is about $2.6$ points above the aggregate model prediction, a gap consistent with worker selection, work intensification, and sectoral capital deepening that the model does not encode. The Brazilian 1988 evidence points in the same direction, with a slightly larger empirical hourly-wage response. These comparisons are directional benchmarks rather than lower bounds for Brazil's effective $A_{\text{req}}$. Portugal and Brazil 1988 differ from the current Brazilian 36h/40h policy environment in informality, institutions, and reform magnitude.

\begin{table}[htbp]
\centering
\caption{Calibration checks and directional benchmarks.}\label{tab:calibration_fit}
\small
\begin{tabular}{lccc}
\hline\hline
Moment & Observed & Model & Gap \\
\hline
\multicolumn{4}{l}{\textit{Panel A: Calibration targets (matched by construction)}} \\
\hline
Informality rate, small firms ($\leq 49$ empl.)  & 0.500 & 0.500 & 0.000 \\
Informality rate, large firms                    & 0.200 & 0.200 & 0.000 \\
Formal employment share, small firms             & 0.590 & 0.590 & 0.000 \\
Formal employment share, large firms             & 0.410 & 0.410 & 0.000 \\
Aggregate informality, narrow PNAD               & 0.378 & 0.377 & 0.001 \\
Wage premium $R$, MNR (2015) central target       & 1.400 & 1.400 & 0.000 \\
\hline
\multicolumn{4}{l}{\textit{Panel B1: Domestic checks outside calibration}} \\
\hline
Average weekly formal hours, habitual measure            & 41.23 & 42.24 & 1.01 h/week \\
Formal share above $36$h, habitual measure (\%)          & 85.7  & 91.5  & 5.80 pp \\
Sectoral range in average formal hours                   & 3.31  & 0.80  & $-2.51$ h/week \\
\hline
\multicolumn{4}{l}{\textit{Panel B2: External directional benchmarks}} \\
\hline
$\Delta Y$ at $44\to 40$h cap, Portugal 1996 (firm-level) & $-3.20$ & $-2.00$ & 1.20 pp \\
$\Delta Y/h$ at $44\to 40$h cap, Portugal 1996 (\%)       & $+4.40$ & $+1.77$ & 2.63 pp \\
$\Delta W/h$ at $48\to 44$h, Brazil 1988 (\%, GMC 2003)   & $+7.00$ & $+1.77$ & 5.23 pp \\
\hline\hline
\end{tabular}
\\[2pt]
\begin{minipage}{0.95\textwidth}
\justifying
{\footnotesize Panel B1 rows 1 and 2 compare PNAD habitual-hours moments with the baseline model moments implied by the contracted-hours calibration. Row 3 reports the range across agriculture, industry, and services in average formal weekly hours. The model value comes from the sectoral extension using the same hours-bin structure, so the negative gap means that the model understates observed sectoral dispersion in formal hours. The aggregate-hours gap is driven mainly by the maintained assumption $h_I=44$ for informal workers and is treated as a sensitivity margin in the online appendix, not as a validation target ($A_{\text{req}}$ shifts by $<0.2$ pp for $h_I\in[38,46]$). Panel B2 compares the model's aggregate prediction at the $44\to 40$h cap to the Portuguese firm-level estimate \citep{AsaiLopesTondini2024} and to Brazil's 1988 reform of comparable cap-cut magnitude (48$\to$44h, $\sim$8\%) as documented by \citet{GonzagaMenezesFilhoCamargo2003}. Both external benchmarks suggest the model is conservative on the hourly-productivity response. The model's $e(h)$ captures fatigue and not worker selection, work intensification, or sectoral capital deepening. Sources: PNAD/SIDRA, RAIS 2022, DIEESE 2024, \citet{MeghirNaritaRobin2015}, \citet{AsaiLopesTondini2024}, and \citet{GonzagaMenezesFilhoCamargo2003} \citep{PNAD2025,RAIS2022,DIEESE2024}.}
\end{minipage}
\end{table}


% ==================================================================
% 4. Main results (includes welfare as Section 4.4)
% ==================================================================
% =============================================================
% Section: Main results
% =============================================================

\section{The productivity arithmetic of alternative hours caps}\label{sec:results}

The main lesson is that a 40-hour cap and a 36-hour cap are different economic objects. Moving from 44 to 40 hours cuts the long tail of the contracted-hours distribution and lets many workers move toward the assumed efficiency peak. Moving below 40 hours is no longer mainly about trimming fatigue. It asks the economy to give up hours that the model no longer treats as inefficient. That is why the curve in Figure~\ref{fig:areq_curve} bends upward rather than sliding down smoothly.

Table~\ref{tab:main_results} reports the 36-hour endpoint, but the useful comparison is the slope of the path. The 40-hour alternative needs only about a $2\%$ one-off productivity gain to restore output. The 36-hour endpoint needs a much larger productivity response: $6.63\%$ under the one-sided fatigue benchmark and $8.18\%$ under the two-sided penalty benchmark. The distinction between these two columns is not a technicality. In the first, shorter workweeks are not assigned a direct efficiency penalty once hours fall below 40. In the second, moving away from 40 in either direction is costly.

Informality moves in the expected direction, but it is not the villain of the aggregate result. The calibrated economy does shift some marginal jobs toward informality, and the increase matters for incidence. Still, most of the output loss comes from the simpler arithmetic of fewer formal hours. The joint $(\sigma,\omega,\eta_I)$ envelope reinforces that point: changing formal-informal substitutability, the formal CES weight, and informal relative efficiency moves the endpoint requirement, but it does not turn the 36-hour reform into a small productivity problem. Capital adjustment can shave roughly one to two percentage points from the short-run benchmark over a three-to-five-year horizon, but it does not erase the difference between 40 and 36 hours.

\begin{table}[htbp]
\centering
\caption{Main results at the 36-hour endpoint.}\label{tab:main_results}
\small
\begin{tabular}{@{}l c c@{}}
\hline\hline
 & One-sided fatigue & Two-sided penalty \\
 & (preferred) & (conservative) \\
\hline
\multicolumn{3}{l}{\textit{Positive predictions (no TFP compensation)}} \\
Output, $\Delta Y$                       & $-6.60\%$        & $-8.02\%$ \\
Welfare (GHH), $\Delta$CV                & $-1.76\%$        & $-3.63\%$ \\
Informality, $\Delta$Inf                 & $+1.57$ pp       & $+1.92$ pp \\
Output per hour, $\Delta Y/h$            & $+2.39\%$        & $+0.75\%$ \\
\hline
\multicolumn{3}{l}{\textit{Benchmark: TFP gain that restores output}} \\
$A_{\text{req}}$, central                & $6.63\%$         & $8.18\%$ \\
$A_{\text{req}}$, joint envelope         & $[5.62,\,7.48]\%$ & $[6.93,\,9.23]\%$ \\
\hline\hline
\end{tabular}
\\[2pt]
\begin{minipage}{0.95\textwidth}
\justifying
{\footnotesize The joint envelope spans the eight corners of the disciplined $(\sigma,\omega,\eta_I)$ box. The one-sided fatigue envelope lies below the two-sided penalty envelope because the former does not penalize efficiency below 40 hours.}
\end{minipage}
\end{table}

\begin{figure}[htbp]
\centering
\includegraphics[width=0.92\textwidth]{fig_areq_vs_hours.pdf}
\caption{Required TFP gain $A_{\text{req}}$ as a function of the weekly hours cap.}\label{fig:areq_curve}
\vspace{2pt}
\begin{minipage}{0.92\textwidth}
\justifying
{\footnotesize \textit{Notes:} The vertical markers distinguish the 40h alternative in the relator proposal from the 36h endpoint in the original PEC 8/2025 proposal. The one-sided fatigue schedule yields $6.63\%$ at the $36$h cap. The two-sided penalty schedule yields $8.18\%$. The two schedules coincide above $h^{*}=40$ and diverge below it, reflecting different assumptions about $e(h)$ below the efficiency peak.}
\end{minipage}
\end{figure}

The figure also clarifies what the model is, and is not, saying. The first four hours of reduction are cheap in productivity units because the hours-efficiency schedule is doing helpful work. The next four hours are expensive because the reform is now mostly removing effective labor. Appendix Table~\ref{tab:efficiency_diagnostics} reports the diagnostic checks behind this statement. With no hours-efficiency channel at all, the 36-hour endpoint would require $7.74\%$ rather than $6.63\%$ or $8.18\%$; shifting the efficiency peak by one hour changes the endpoint more than changing several secondary parameters. This is why I treat the 41--42h welfare maximum below as a model result, not as a free-standing policy rule.

The historical TFP comparison should be read in the same spirit. PWT 11.0 is not a validation exercise; it is a ruler. Brazil's post-1990 TFP record is weak, so a one-off gain of six to eight percent is a large number on that ruler. The 40-hour alternative lives on a different part of the scale. Appendix Table~\ref{tab:tfp_history} translates the one-off requirements into annualized gains over three-, five-, and ten-year phase-ins. That phase-in framing is more useful than asking whether a single historical decade mechanically ``matches'' the counterfactual.

\begin{figure}[htbp]
\centering
\includegraphics[width=0.95\textwidth]{fig_areq_vs_tfp_history.pdf}
\caption{Brazilian TFP and the required one-off productivity gain.}\label{fig:tfp_history}
\vspace{2pt}
\begin{minipage}{0.95\textwidth}
\justifying
{\footnotesize \textit{Notes:} Panel~A plots the PWT 11.0 \texttt{rtfpna} series (real TFP at constant national prices) for Brazil, 1954 to 2023. Panel~B shows annualized TFP growth rates by horizon. Labels indicate the years required to accumulate the preferred 36-hour baseline $A_{\text{req}}=6.63\%$ at each historical rate (``never'' when the rate is non-positive). \textit{Source:} Penn World Table 11.0 via FRED series \texttt{RTFPNABRA632NRUG} \citep{FeenstraInklaarTimmer2015,PWT2025,FREDPWTBrazil2026}.}
\end{minipage}
\end{figure}

A downside stress test in the appendix asks what happens if the reorganization itself temporarily lowers productivity. This is not the central prediction, but it is a useful guardrail: a one-percent post-reform TFP disruption pushes the two-sided endpoint requirement above nine percent. The reason to report this case is not to forecast disruption. It is to show that the policy arithmetic is sensitive to implementation quality, not only to steady-state labor supply.

% =============================================================
% Subsection: Welfare and heterogeneity (inside Section 4)
% =============================================================

\subsection{Welfare and incidence}\label{sec:welfare}

The welfare calculation is best read as a consumption-leisure diagnostic. It asks a narrow question: if employment is fixed, wages do not bargain, and the reform mechanically gives incumbent workers more non-work time, how large is the leisure gain relative to the consumption loss? It is not a social-welfare theorem. Health, household production, bargaining, prices, employment transitions, and distributional weights are outside the calculation.

Within that narrow object, the schedule has a clear shape. Around 41 to 42 hours, the model buys leisure at a relatively low output cost. Below 40 hours, the price of additional leisure rises sharply because the reform is no longer mostly trimming long-hours fatigue. That is why the representative-agent welfare measure turns negative before the endpoint. The one-sided fatigue and two-sided penalty specifications agree on the moderate-reduction region and disagree mainly once the cap moves below the efficiency peak.

\begin{figure}[htbp]
\centering
\includegraphics[width=0.92\textwidth]{fig_welfare_schedule.pdf}
\caption{Welfare schedule by hours cap under both calibrations.}\label{fig:welfare_schedule}
\vspace{2pt}
\begin{minipage}{0.92\textwidth}
\justifying
{\footnotesize \textit{Notes:} Each line plots aggregate welfare $\Delta$CV from Equation~\eqref{eq:cv} at successively lower hours caps, relative to the $44$h baseline, under the one-sided fatigue and two-sided penalty $e(h)$ specifications. $\Delta$CV peaks near $42$h at about $+0.62\%$ and turns negative below $40$h. At the $36$h endpoint, welfare falls by $1.76\%$ under the one-sided specification and by $3.63\%$ under the two-sided one.}
\end{minipage}
\end{figure}

A TFP gain of approximately $1.4\%$ at the 36h cap is enough to make the endpoint welfare-neutral under the one-sided fatigue specification. The two-sided threshold is $2.9\%$. Both thresholds lie below the output-neutral $A_{\text{req}}$ because leisure partially compensates the consumption loss. Appendix Table~\ref{tab:welfare_preference_robustness} shows that this qualitative ranking survives a separable log and CRRA(2) diagnostic: the 40-hour alternative remains positive, while the 36-hour endpoint remains negative.

The aggregate welfare number masks incidence. In the conservative type-level decomposition reported in Appendix Table~\ref{tab:welfare_incidence_by_type}, formal stayers gain $+1.41\%$ in first-order compensating variation because they work roughly six fewer hours. Workers pushed into informality lose $-28.57\%$. This group is small, $1.9\%$ of workers, but its loss is large. Informal workers not affected by the statutory cap have zero direct incidence in the model. The representative-agent value in Equation~\eqref{eq:cv} differs from the share-weighted type sum because the GHH composite is nonlinear. For this reason, the 41 to 42h welfare maximum should be read as a model-internal representative-agent result, with distributional incidence reported separately.

The welfare schedule does not imply opposition to reductions in working hours. Under the model's calibration, the welfare-improving margin lies near 41 to 42 hours rather than at 36. The schedule separates output loss, required productivity, and welfare. These objects move differently with the cap, and the welfare-improving region is narrower than the output-neutral one.

Adjustment costs and informality differ across firm sizes and sectors, producing heterogeneous incidence of the reform. Small firms ($\leq 49$ employees) require a larger TFP gain than large firms, reflecting their higher formalization wedges and baseline informality. Applying the national framework separately to three broad sectors (agriculture, industry, services) with sector-specific PNAD parameters delivers a per-sector $A_{\text{req}}$ of $1.3$ to $1.8\%$ for the 40-hour alternative and $6.91$ to $8.05\%$ for the 36-hour endpoint under the two-sided penalty specification. Informality rises in every sector at the 36-hour endpoint (agriculture $+2.88$ pp, industry $+2.73$ pp, services $+2.10$ pp). These are sector-by-sector partial-equilibrium counterfactuals, not an input-output incidence analysis. Industry's $A_{\text{req}}$ slightly exceeds agriculture's because the moderate $\eta_I=0.40$ combined with industry's larger formal base dampens the mechanical gap between the two sectors. Services account for most of the aggregate output loss through their output share, not through a larger per-worker effect. The online appendix reports the full decomposition and the parameter sensitivities.


% ==================================================================
% 5. Policy implications
% ==================================================================
% =============================================================
% Section: Policy implications
% =============================================================

\section{Implications for reform design}\label{sec:conclusion}

The relevant policy object is no longer only the original 36-hour version of PEC 8/2025. The live debate also includes 40-hour, 5x2, and phased alternatives. The model therefore reports the full hours-cap schedule. A 40-hour cap has a much smaller output-neutrality benchmark than the endpoint and leaves the representative agent GHH diagnostic positive. A single-step move from 44 to 36 hours is different. It affects most Brazilian formal workers and reduces the labor input of the economy by about $14\%$ mechanically. At that endpoint, the output loss is offset by a one-off productivity gain of $6.63\%$ under one-sided fatigue and $8.18\%$ under the two-sided penalty.

The 36-hour endpoint is large relative to Brazil's recent productivity record. The PWT comparison should be read as a scale benchmark, because an annual macro residual is not the same object as the short-run productivity response in the model. It is nevertheless informative about magnitudes: post-1990 Brazilian TFP growth has been weak, so the endpoint would require an unusual acceleration if output were to be preserved over a short implementation horizon. Without such a productivity response, the model implies a short-run output loss of $6.6$ to $8.0\%$ and a decline in the aggregate GHH diagnostic of $1.8$ to $3.6\%$ at the endpoint.

A moderate reduction can raise the model's GHH diagnostic. In the representative agent calibration, cutting the workweek from $44$ to $41$ or $42$ hours raises that indicator because the value assigned to extra leisure exceeds the consumption loss. The 40-hour alternative is near the edge of the positive region, and below $40$h the calculation turns negative. This result depends on the assumed hours-productivity schedule and representative preferences. It is not a social welfare ordering over heterogeneous workers.

The cost of the endpoint is concentrated. Small firms bear the largest per-worker productivity requirement. Industry and agriculture carry the largest per-sector requirements, and informality rises in all three sectors. Services account for most of the aggregate output loss through their employment and output share, not through a larger per-worker effect. The sectoral exercise is partial-equilibrium and sector-by-sector rather than an input-output incidence model, but it is enough to show that the adjustment burden is uneven.

Informality rises but does not dominate the output cost. The calibrated model, disciplined by the observed formal-informal wage premium, finds an aggregate informality increase of $1.6$ to $1.9$ pp. That rise matters for welfare accounting, especially for workers pushed into informality. It is small relative to the mechanical hours contraction. The output cost comes primarily from fewer hours worked.

These conclusions have limits. The model is short-run and partial-equilibrium. Endogenous capital adjustment over a three-to-five-year horizon lowers the required productivity gain by about one to two percentage points. Total employment is fixed, which shuts down work-sharing. Evidence on work-sharing from Germany \citep{Hunt1996} and Brazil's 1988 constitutional reform \citep{GonzagaMenezesFilhoCamargo2003} suggests that this restriction is compatible with available evidence, but it is still a restriction. Compliance, overtime, organizational disruption, wage bargaining, and informality responses could move the relevant benchmark in either direction. The Portuguese 1996 experience \citep{AsaiLopesTondini2024} documents firm-level productivity gains larger than the model produces, which suggests that the model is conservative on the hourly-productivity response.

Under these limits, the productivity arithmetic suggests that phased paths deserve separate evaluation. An intermediate stop near $40$ to $42$ hours keeps the productivity requirement far below the 36-hour endpoint and gives capital, informality, and firm-level productivity time to adjust before a further step. The implication is not that the workweek should remain at 44 hours, nor that the model identifies the socially optimal cap. It is that the size and timing of the reduction matter for the tradeoff between output neutrality, representative agent welfare, and incidence across firms and workers.


% === Online appendix and replication ===
\paragraph{Online appendix.} The \href{\onlineappendixurl}{public replication repository} contains the online appendix with data sources, omitted-channel diagnostics, calibration details, sensitivity grids, sectoral results, welfare incidence, output decomposition, historical TFP table, derivations, and hours-efficiency robustness.

\paragraph{Code and data availability.} Code, calibration targets, and unit tests to reproduce every numbered table and figure in this paper are available at the author's \href{https://github.com/vr-rodrigues/jornada-6x1-replication}{public GitHub repository}.

% === Bibliography ===
\bibliographystyle{aer}
\bibliography{bibliography_clean}

\end{document}
